% ============================================================
%  Chapter 15: Triple Equivalence  𝔒 ≅ 𝔈 ≅ 𝔓
%  Source: union/publication/support/unconstrained-subtask-recursion.tex
%  Parts II, VII, IX--X
% ============================================================

\epigraph{%
  Classical mechanics and quantum mechanics\\
  are not two different theories.\\
  They are two functors from the same category.%
}{}

%% ── §15.1 ─────────────────────────────────────────────────────────────────────
\section{Three Categories, One Bounded Phase Space}
\label{sec:ch15-three}

Every state of a physical system admits three equivalent descriptions,
each corresponding to a distinct observational language:

\begin{enumerate}[(i)]
  \item The \textbf{Oscillatory category} $\Osc$: objects are oscillatory
    structures $(M, \omega, \Phi, \mu)$ (Definition~\ref{def:osc-mono});
    morphisms are frequency- and phase-preserving maps.  Language: modes,
    frequencies, Fourier representation.  This is the natural language of
    classical mechanics.

  \item The \textbf{Categorical (Epistemological) category} $\Cat$: objects are
    categorical structures $(C, \prec, \lambda, \nu)$
    (Definition~\ref{def:cat-mono}); morphisms preserve the refinement order
    and labelling.  Language: S-entropy, refinement hierarchies,
    epistemological distance.  This is the language of the S-entropy calculus.

  \item The \textbf{Partition category} $\Part$: objects are partition
    structures $(\Omega, \mathscr{P}, \pi)$ (Definition~\ref{def:part-mono});
    morphisms preserve refinement and the measure.  Language: partition cells,
    depths, measure-theoretic structure.  This is the natural language of
    quantum mechanics (cells correspond to eigenstates; refinement maps
    correspond to quantum channels).
\end{enumerate}

The central theorem of this chapter---the Triple Equivalence---asserts that
these three categories are equivalent as categories.  The equivalence is not
approximate or asymptotic; it holds exactly, via explicit conversion functors.

%% ── §15.2 ─────────────────────────────────────────────────────────────────────
\section{The Three Conversion Functors}
\label{sec:ch15-functors}

We construct the three conversion functors explicitly before stating the
equivalence theorem.

\begin{definition}[Functor $F_{OC}$: Oscillatory to Categorical]
\label{def:F-OC}
Define $F_{OC} : \mathbf{Osc} \to \mathbf{Cat}$ on objects by
\[
  F_{OC}(M, \omega, \Phi, \mu)
  \;:=\; (M,\; \prec_\omega,\; \lambda_\omega,\; \mu)
\]
where:
\begin{itemize}[noitemsep]
  \item $m_1 \prec_\omega m_2$ if and only if $\omega(m_1)$ divides
    $\omega(m_2)$ in $\mathbb{R}_{>0}$ (interpreted as a rational ratio with
    bounded denominator), defining the refinement order from frequency
    divisibility;
  \item the labelling
    $\lambda_\omega(m) :=
    \bigl(\lfloor\log_2 \omega(m)\rfloor,\;
     \lfloor(\omega(m) - 2^{\lfloor\log_2\omega(m)\rfloor})/
            2^{\lfloor\log_2\omega(m)\rfloor - 1}\rfloor,\;
     \arg\Phi(m,0)\bigr)$
    encodes the octave, the sub-octave position, and the initial phase of
    each mode.
\end{itemize}
On morphisms, $F_{OC}$ acts by restriction.
\end{definition}

\begin{definition}[Functor $F_{CP}$: Categorical to Partition]
\label{def:F-CP}
Define $F_{CP} : \mathbf{Cat} \to \mathbf{Part}$ on objects as follows.
Given $(C, \prec, \lambda, \nu)$, let $\Omega := C / {\sim_n}$ where
${\sim_n}$ identifies cells whose labels agree on the first $n$ coordinates;
the family $\{P_n\}$ of partitions is induced by this identification.
Define $\pi$ as the pushforward of $\nu$ under the quotient by ${\sim_n}$ in
the limit $n \to \infty$.
\end{definition}

\begin{definition}[Functor $F_{PO}$: Partition to Oscillatory]
\label{def:F-PO}
Define $F_{PO} : \mathbf{Part} \to \mathbf{Osc}$ by
\[
  F_{PO}(\Omega, \mathscr{P}, \pi)
  \;:=\; (P_\infty,\; \omega_{\mathfrak{P}},\; \Phi_{\mathfrak{P}},\; \pi)
\]
where:
\begin{itemize}[noitemsep]
  \item $P_\infty = \bigvee_n P_n$ is the limit partition (the atoms of the
    generated $\sigma$-algebra);
  \item $\omega_{\mathfrak{P}}(p) := 1/\pi(p)$ is the inverse-measure
    frequency, by analogy with Kac's lemma (the mean return time to a set
    of measure $\pi(p)$ is $1/\pi(p)$);
  \item $\Phi_{\mathfrak{P}}(p, t) := e^{2\pi i\,\omega_{\mathfrak{P}}(p)\,t}$.
\end{itemize}
\end{definition}

%% ── §15.3 ─────────────────────────────────────────────────────────────────────
\section{The Triple Equivalence Theorem}
\label{sec:ch15-equivalence}

\begin{theorem}[Triple Equivalence]\label{thm:triple-equiv-mono}
The three categories $\mathbf{Osc}$, $\mathbf{Cat}$, $\mathbf{Part}$ are
equivalent as categories.  Specifically, the functors $F_{OC}$, $F_{CP}$,
$F_{PO}$ are part of mutually inverse equivalences:
\begin{align}
  F_{CO} \circ F_{OC} &\cong \mathrm{id}_{\mathbf{Osc}}, &
  F_{OC} \circ F_{CO} &\cong \mathrm{id}_{\mathbf{Cat}}, \label{eq:OC-eq}\\
  F_{PC} \circ F_{CP} &\cong \mathrm{id}_{\mathbf{Cat}}, &
  F_{CP} \circ F_{PC} &\cong \mathrm{id}_{\mathbf{Part}}, \label{eq:CP-eq}\\
  F_{OP} \circ F_{PO} &\cong \mathrm{id}_{\mathbf{Part}}, &
  F_{PO} \circ F_{OP} &\cong \mathrm{id}_{\mathbf{Osc}}. \label{eq:PO-eq}
\end{align}
Moreover, composing any two of the three forward functors yields the third up
to natural isomorphism:
\[
  F_{OC} \cong F_{PC} \circ F_{OP}, \quad
  F_{CP} \cong F_{OP} \circ F_{CO}, \quad
  F_{PO} \cong F_{CO} \circ F_{PC}.
\]
\end{theorem}

\begin{proof}
We construct the inverse functor $F_{CO} : \mathbf{Cat} \to \mathbf{Osc}$;
the inverses $F_{PC}$ and $F_{OP}$ follow by symmetric arguments.

Given $\mathfrak{E} = (C, \prec, \lambda, \nu)$, define
$F_{CO}(\mathfrak{E}) := (M, \omega, \Phi, \mu)$ with $M := C$,
$\omega(c) := \exp(\lambda_1(c) \ln 2)$ (recovering the frequency from the
first label coordinate, the octave),
$\Phi(c, 0) := e^{i\lambda_3(c)}$ (recovering the initial phase from the
third label coordinate), extended to all $t$ by the phase law of
Definition~\ref{def:osc-mono}, and $\mu := \nu$.

The composition $F_{CO} \circ F_{OC}$ acts on $(M, \omega, \Phi, \mu)$ by
computing the labels via $F_{OC}$ and then inverting via $F_{CO}$; up to the
rounding ambiguity in $\lfloor\cdot\rfloor$, the result is isomorphic to the
identity (equivalence, not equality, of categories).

The naturality and triangle identities are verified by direct computation on
morphisms.  The key observation is that all three structures encode the same
information about the underlying probability measure under the same generating
$\sigma$-algebra; the choice between them is computational, not informational.
\end{proof}

%% ── §15.4 ─────────────────────────────────────────────────────────────────────
\section{Commutative Diagram}
\label{sec:ch15-diagram}

The Triple Equivalence has the structure of a commuting triangle of
category equivalences.  The following diagram illustrates the six functors
(three forward, three inverse) and the cyclic structure:

\bigskip
\begin{center}
\begin{tikzpicture}[
  every node/.style={font=\normalsize},
  node distance=4.5cm,
  >=stealth,
  thick
]
  % Nodes
  \node (Osc)  at (0,0)    {$\mathbf{Osc}\;\bigl(\Osc\bigr)$};
  \node (Cat)  at (6,0)    {$\mathbf{Cat}\;\bigl(\Cat\bigr)$};
  \node (Part) at (3,-4.5) {$\mathbf{Part}\;\bigl(\Part\bigr)$};

  % Forward functors (outer arcs, curving away from centre)
  \draw[->] (Osc)  to[bend left=15]
    node[above, midway] {$F_{OC}$} (Cat);
  \draw[->] (Cat)  to[bend left=15]
    node[right, midway, xshift=3pt] {$F_{CP}$} (Part);
  \draw[->] (Part) to[bend left=15]
    node[left, midway, xshift=-3pt] {$F_{PO}$} (Osc);

  % Inverse functors (inner arcs, curving toward centre)
  \draw[->, dashed] (Cat)  to[bend left=15]
    node[below, midway] {$F_{CO}$} (Osc);
  \draw[->, dashed] (Part) to[bend left=15]
    node[left,  midway, xshift=3pt] {$F_{PC}$} (Cat);
  \draw[->, dashed] (Osc)  to[bend left=15]
    node[right, midway, xshift=-3pt] {$F_{OP}$} (Part);

  % Equivalence labels at vertices
  \node[below=0.08cm of Osc,  font=\small, color=gray]
    {classical / oscillatory};
  \node[below=0.08cm of Cat,  font=\small, color=gray]
    {epistemological / S-entropy};
  \node[below=0.08cm of Part, font=\small, color=gray]
    {partition / quantum};
\end{tikzpicture}
\end{center}
\bigskip

\noindent
Solid arrows are the forward functors; dashed arrows are the inverse functors.
The diagram commutes up to natural isomorphism in every direction.

%% ── §15.5 ─────────────────────────────────────────────────────────────────────
\section{Free Conversion and Optimal Representation}
\label{sec:ch15-free-conv}

\begin{corollary}[Free Conversion]\label{cor:free-conv-mono}
For any computation $f$ defined in one of the three representations, there
exist computations $f^O, f^C, f^P$ in the other representations such that the
diagram of conversions commutes up to natural isomorphism.  A computation may
therefore be performed in whichever representation makes it cheapest.
\end{corollary}

\begin{proof}
Direct from Theorem~\ref{thm:triple-equiv-mono}: the equivalence functors
translate any computation between representations without information loss.
\end{proof}

\begin{definition}[Computation cost and conversion cost]
\label{def:costs-mono}
For a computation $f$ in representation $R \in \{\mathfrak{O}, \mathfrak{E}, \mathfrak{P}\}$,
the \emph{computation cost} $\mathrm{Conv}_R(f) \in \mathbb{R}_{\geq 0}$ is
the minimum number of elementary operations of $R$'s algebraic structure.
The \emph{conversion cost} $\mathrm{Conv}_{R \to R'} \in \mathbb{R}_{\geq 0}$
is the minimum number of elementary operations to convert an object from
representation $R$ to $R'$.
\end{definition}

\begin{theorem}[Optimal Representation]\label{thm:opt-rep-mono}
For any computation $f$, the optimal total cost over the choice of
representation is
\[
  \mathrm{Conv}^*(f)
  = \min_{R \in \{\mathfrak{O},\mathfrak{E},\mathfrak{P}\}}
    \Bigl[\mathrm{Conv}_R(f)
          + \min_{R'}\mathrm{Conv}_{R \to R'}\Bigr].
\]
\end{theorem}

\begin{proof}
Direct from Corollary~\ref{cor:free-conv-mono}: every computation has
equivalent forms in all three representations, so the cost is the minimum of
the representation-specific costs plus any necessary conversion to the desired
output representation.
\end{proof}

\begin{remark}[The apples-and-oranges resolution]
Theorem~\ref{thm:opt-rep-mono} formalises the principle that quantities of
nominally different types (e.g., position and time) can be combined freely
\emph{after} conversion to a common representation.  No special mixed-type
operation is required; the apparent mixing is a consequence of selecting the
most computationally efficient representation in which the calculation is
type-uniform.  This is the formal resolution of the apples-and-oranges problem
that the Unconstrained Subtask Theorem (Chapter~\ref{ch:floor}) addresses at
the level of S-expressions.
\end{remark}

%% ── §15.6 ─────────────────────────────────────────────────────────────────────
\section{The Structural Symmetry Group}
\label{sec:ch15-symmetry}

\begin{definition}[Structural automorphism]\label{def:struct-auto-mono}
A \emph{structural automorphism} of the calculus is a triple of compatible
self-equivalences $(\sigma_O, \sigma_C, \sigma_P)$ on the three representation
categories such that
\[
  \sigma_C \circ F_{OC} = F_{OC} \circ \sigma_O, \quad
  \sigma_P \circ F_{CP} = F_{CP} \circ \sigma_C, \quad
  \sigma_O \circ F_{PO} = F_{PO} \circ \sigma_P.
\]
\end{definition}

\begin{theorem}[Symmetry Group]\label{thm:sym-group-mono}
The structural automorphisms form a group $G_{\mathrm{str}}$ acting on the
algebra of S-expressions $\mathcal{E}$.  The group includes:
\begin{enumerate}[(i)]
  \item the trivial action (identity);
  \item the cyclic permutation $\mathfrak{O} \to \mathfrak{E} \to
    \mathfrak{P} \to \mathfrak{O}$;
  \item the involutive swap of any two representations;
  \item all compositions of (ii) and (iii).
\end{enumerate}
The group $G_{\mathrm{str}}$ is isomorphic to the symmetric group $S_3$.
\end{theorem}

\begin{proof}
Direct verification that compositions of these operations remain structural
automorphisms by Theorem~\ref{thm:triple-equiv-mono}.  The six elements of
$S_3$ correspond to the six ways of permuting three representations.
\end{proof}

\begin{corollary}[Invariants of $G_{\mathrm{str}}$]\label{cor:invariants-mono}
The S-value $S_{\mathcal{R}}$, the catalytic power $\kappa$, and the floor
$S_\flat$ are all invariants of the structural automorphism group: they take
the same value on all symmetric copies of an expression.
\end{corollary}

\begin{proof}
Each of these quantities is defined via the receiver evaluation
$\mathrm{eval}_{\mathcal{R}}$, which commutes with the conversion functors by
definition of the equivalence.  Hence permuting representations leaves S-values,
catalytic powers, and floors unchanged.
\end{proof}

%% ── §15.7 ─────────────────────────────────────────────────────────────────────
\section{Triple Entropy Equivalence}
\label{sec:ch15-entropy}

The Triple Equivalence of categories yields a parallel equivalence of the three
S-entropy components of the mass spectrometric triple.

\begin{consequence}{Triple Entropy Equivalence}{tripleent}
For any fixed partition depth $n$ and any state in the bounded phase space
$(\Omega, \mathscr{P}, \pi)$, the three S-entropy components of the triple
$(\Sk, \St, \Se)$ are equal:
\[
  \Sk \;=\; \St \;=\; \Se.
\]
That is: the oscillatory S-entropy, the categorical S-entropy, and the
partition S-entropy are the same quantity computed in three distinct
observational languages.
\end{consequence}

\begin{proof}
We show that all three S-values coincide with the partition entropy at depth
$n$.

\medskip\noindent
\textit{Partition entropy $\Se$.}
By Definition~\ref{def:part-mono}, the partition structure at depth $n$ has
cells $\{P_n\}$ with weights $\pi(p)$.  The partition S-entropy is
\[
  \Se = -\sum_{p \in P_n} \pi(p)\log\pi(p),
\]
the Shannon entropy of the partition measure at depth $n$.

\medskip\noindent
\textit{Categorical entropy $\St$.}
The functor $F_{PC} : \mathbf{Part} \to \mathbf{Cat}$ assigns to each
partition atom $p \in P_n$ a cell $c_p \in C$ with weight $\nu(c_p) = \pi(p)$.
The categorical S-entropy is
\[
  \St = -\sum_{c \in C_n} \nu(c)\log\nu(c)
       = -\sum_{p \in P_n} \pi(p)\log\pi(p)
       = \Se,
\]
since $F_{PC}$ preserves the measure by construction.

\medskip\noindent
\textit{Oscillatory entropy $\Sk$.}
The functor $F_{PO} : \mathbf{Part} \to \mathbf{Osc}$ assigns to each atom
$p \in P_\infty$ the frequency $\omega(p) = 1/\pi(p)$ and weight $\mu(p) =
\pi(p)$.  The oscillatory S-entropy, computed as the weighted Shannon entropy
of the mode measure, is
\[
  \Sk = -\sum_{p \in P_n} \mu(p)\log\mu(p)
       = -\sum_{p \in P_n} \pi(p)\log\pi(p)
       = \Se.
\]
Therefore $\Sk = \St = \Se$ at any fixed partition depth.
\end{proof}

\begin{remark}[What the equality means]
The equality $\Sk = \St = \Se$ is not a numerical coincidence: it is the
quantitative expression of the Triple Equivalence.  The three categories encode
the same probability measure; the three entropies are three formulas for the
same quantity applied to the same measure.  Choosing which formula to compute
is a matter of computational convenience, not a choice about which entropy is
``real''.
\end{remark}

%% ── §15.8 ─────────────────────────────────────────────────────────────────────
\section{Physical Consequences of the Triple Equivalence}
\label{sec:ch15-consequences}

The Triple Equivalence has three immediate physical consequences, each of
which appears startling from a conventional perspective but follows directly
from Theorem~\ref{thm:triple-equiv-mono}.

%% Consequence 1
\begin{consequence}{Classical Mechanics $=$ Quantum Mechanics}{classeqquant}
Every trajectory of a bounded oscillator (object of $\mathbf{Osc}$, the
language of classical mechanics) corresponds, via $F_{OC} \circ F_{CP}$, to
a partition atom of the limit partition (object of $\mathbf{Part}$, the
language of quantum mechanics), and vice versa via $F_{PO}$.  The two
descriptions are not approximate; they are equivalent.
\end{consequence}

\begin{proof}
Apply Theorem~\ref{thm:triple-equiv-mono}: the functor $F_{OC}$ sends each
classical mode $(m, \omega(m), \Phi(m,\cdot), \mu)$ to a categorical cell, and
$F_{CP}$ sends that cell to a partition atom.  The inverse $F_{PO} \circ F_{CO}$
reconstructs the classical mode from the partition atom.  The reconstruction
is exact up to the rounding ambiguity of the categorical labelling, which
vanishes in the equivalence of categories.
\end{proof}

\begin{remark}[Scope of the consequence]
The statement ``classical mechanics $=$ quantum mechanics'' is scoped
to bounded physical systems satisfying Axiom~\ref{ax:bpsl}.  For such
systems, the Koopman operator on $L^2(\Omega, \mu)$ is unitary (by
clause~(ii) of the Axiom), and its eigenfunctions are the oscillatory modes.
The partition atoms are the cells of the hierarchical partition guaranteed by
clause~(iii).  The Triple Equivalence says that the eigenfunction description
and the partition-cell description carry identical information.  This does not
imply that classical mechanics and quantum mechanics make the same predictions
in all regimes; it implies that, for persistent bounded systems, the distinction
is one of computational language, not of physical content.
\end{remark}

%% Consequence 2
\begin{consequence}{GPU Computation $=$ Physical Instrument}{gpuphys}
A GPU fragment shader performing partition-depth minimisation on a bounded
domain (an object of $\mathbf{Part}$) is a physical observation apparatus in
the sense of the oscillatory category $\mathbf{Osc}$ (via $F_{PO}$) and an
epistemological receiver in the sense of the categorical category $\mathbf{Cat}$
(via $F_{PC}$).  The GPU and the physical instrument compute the same S-value
and attain the same floor.
\end{consequence}

\begin{proof}
Let the GPU's partition be $(\Omega_{\mathrm{GPU}}, \mathscr{P}_{\mathrm{GPU}},
\pi_{\mathrm{GPU}})$ where $\Omega_{\mathrm{GPU}}$ is the shader's bounded
floating-point domain, $\mathscr{P}_{\mathrm{GPU}}$ is the hierarchy of GPU
pixel grids at successive resolutions, and $\pi_{\mathrm{GPU}}$ is the uniform
measure on pixels.  By Theorem~\ref{thm:triple-equiv-mono}, this partition
structure is equivalent to an oscillatory structure via $F_{PO}$ and to a
categorical structure via $F_{PC}$.  The S-value computed by the GPU
(partition entropy at the finest resolution) equals $\Sk = \Se$ by the Triple
Entropy Equivalence (Consequence~\ref{cons:tripleent}).  The floor of the GPU
as a bounded receiver equals $S_\flat(\mathcal{R}_{\mathrm{GPU}}) > 0$ by
Theorem~\ref{thm:floor-pos-mono}, which is the same floor that governs the
physical instrument's mass accuracy.  Hence the GPU shader is, in the sense
of the calculus, a physical observation apparatus.
\end{proof}

%% Consequence 3
\begin{consequence}{S-entropy Uniqueness Across Representations}{sentunique}
For any bounded physical system satisfying Axiom~\ref{ax:bpsl}, the S-entropy
$S_{\mathcal{R}}(x, x^{*})$ is the same number regardless of whether the
receiver computes it in the oscillatory, categorical, or partition
representation.  That is: $S_k = S_t = S_e$ where the subscripts denote the
representation used for the computation.
\end{consequence}

\begin{proof}
This is the quantitative content of Corollary~\ref{cor:invariants-mono}: S-values
are invariants of the structural automorphism group $S_3$ acting on the three
representations.  Therefore $S_{\mathcal{R}}$ takes the same value regardless
of which representation the receiver uses for computation.
\end{proof}

%% ── §15.9 ─────────────────────────────────────────────────────────────────────
\section{Connection to Standard Mathematics}
\label{sec:ch15-math}

The Triple Equivalence connects to several well-established bodies of
mathematics.

\begin{remark}[Stone-type duality]
The equivalence $\mathbf{Cat} \cong \mathbf{Part}$ has the structural shape of
Stone's duality between Boolean algebras and totally disconnected compact
Hausdorff spaces: the categorical structure plays the role of an algebraic
object whose ``points'' correspond to maximal filters, and the partition
structure plays the role of the dual topological space.  The functors $F_{CP}$
and $F_{PC}$ are the Stone-\v{C}ech-style maps in this analogy.
\end{remark}

\begin{remark}[Koopman operator and oscillatory structure]
The equivalence $\mathbf{Part} \cong \mathbf{Osc}$ is the category-theoretic
formulation of the Koopman operator approach to ergodic theory: oscillatory
modes correspond to eigenfunctions of the Koopman operator
$U_t f = f \circ \phi_t$ on $L^2(\Omega, \mu)$, and partition atoms correspond
to the supports of ergodic measures.  The functor $F_{PO}$ is the operator that
converts a partition measure into the corresponding Koopman eigenspectrum.
\end{remark}

\begin{remark}[Ergodic theory and S-convergence]
The catalyst-driven convergence theorem (Theorem~\ref{thm:cat-conv-mono} of
Chapter~\ref{ch:floor}) is structurally similar to Birkhoff's ergodic averaging
theorem: a sequence of operations applied to an initial state yields convergence
to an invariant set, with the floor playing the role of the invariant measure's
support.  The mean return time to a partition cell of measure $\pi(p)$ is
$1/\pi(p)$ by Kac's lemma, which is exactly the frequency $\omega(p)$ assigned
by $F_{PO}$.
\end{remark}

%% ── §15.10 ────────────────────────────────────────────────────────────────────
\section{Cross-Representation Computation}
\label{sec:ch15-cross-rep}

\begin{theorem}[Cross-Representation Composition]\label{thm:cross-rep-mono}
Let $\xi_1, \ldots, \xi_n \in \mathcal{E}$ be expressions in possibly distinct
representations $R_1, \ldots, R_n \in \{\mathfrak{O}, \mathfrak{E}, \mathfrak{P}\}$,
and let $\star_1, \ldots, \star_{n-1}$ be binary operations.  Then the
composition $\xi_1 \star_1 \xi_2 \star_2 \cdots \star_{n-1} \xi_n$ is
well-defined under $\mathrm{eval}_{\mathcal{R}}$ for any
representation-coherent receiver $\mathcal{R}$.
\end{theorem}

\begin{proof}
By the definition of receiver evaluation, the receiver applies the conversion
functor $\Phi^{R \to R'}$ implicitly when an operation requires uniform
representation.  Specifically, when $\star_i$ requires arguments in
representation $R$ but $\xi_i$ is in $R_i \neq R$, the implicit conversion
$\Phi^{R_i \to R}(\xi_i)$ is performed.  The result is well-defined by
Theorem~\ref{thm:triple-equiv-mono} and the conversion cost is bounded by
Definition~\ref{def:costs-mono}.
\end{proof}

\begin{corollary}[Apples-and-Oranges Admissibility]\label{cor:apples-mono}
Let $\mathcal{R}$ be a representation-coherent receiver.  For any finite
sequence of expressions $\xi_i$ each in any representation, and any binary
operation sequence, the composition is admissible at $\mathcal{R}$.
\end{corollary}

\begin{proof}
Apply Theorem~\ref{thm:cross-rep-mono} repeatedly.
\end{proof}

\begin{example}[Three equivalent solutions to the same problem]
Let $x^{*}$ correspond to an oscillatory mode $m^{*}$ with frequency
$\omega(m^{*}) = 2\pi/T$.  Three ostensibly distinct expressions for the
same answer are:
\begin{enumerate}[(1)]
  \item \emph{Oscillatory}: $\xi_O := (m^{*}, \mathfrak{O})$, directly
    identifying the mode by its frequency.
  \item \emph{Categorical}: $\xi_C := (c^{*}, \mathfrak{E})$ where $c^{*}$
    is the cell whose label corresponds to $m^{*}$ under $F_{OC}$.
  \item \emph{Partition}: $\xi_P := (p^{*}, \mathfrak{P})$ where $p^{*}$
    is the partition atom of measure $\pi(p^{*}) = 1/\omega(m^{*})$
    under $F_{PO}$.
\end{enumerate}
By Theorem~\ref{thm:triple-equiv-mono}, all three evaluate to the same
internal state of the receiver and hence have identical S-value.  The choice
is solely a matter of computational cost (Theorem~\ref{thm:opt-rep-mono}).
\end{example}

%% ── §15.11 ────────────────────────────────────────────────────────────────────
\section{The Union: Why Two Crowns}
\label{sec:ch15-union}

The title of this monograph---\emph{Union of Two Crowns}---refers to the
Triple Equivalence as a unifying theorem.  The ``two crowns'' are the two
languages that physics has historically treated as separate: the classical
(oscillatory, continuous, trajectory-based) and the quantum (partition-based,
discrete, eigenstate-based).  The Triple Equivalence shows that both crowns are
functorially equivalent to a third---the epistemological language of S-entropy
and receivers---and therefore to each other.

The union is not a merger by approximation.  It is an exact equivalence of
categories, guaranteed by the explicit construction of the functors $F_{OC}$,
$F_{CP}$, $F_{PO}$ and their inverses.  Any statement made in one language has
an exact translation into the other two; any computation performed in one
representation has a formally equivalent computation in the others.

For mass spectrometry this has immediate content.  The classical description
of an ion trajectory (oscillatory language) and the quantum description of the
ion's energy levels (partition language) are two computations of the same
physical system, related by the conversion functors.  The S-entropy triple
$(\Sk, \St, \Se)$ is a single scalar quantity measured three ways; by the
Triple Entropy Equivalence they are identical.  The GPU implementation of the
mass spectrometric receiver is, in the precise sense of the calculus, a
physical instrument---not a simulation of one.

%% ── Chapter Summary ───────────────────────────────────────────────────────────
\bigskip
\begin{tcolorbox}[colback=crowngold!10, colframe=crowngold!70!black,
                  title={\textbf{Chapter 15 Summary: The Union}}]
\begin{itemize}[noitemsep]
  \item Three categories are defined: $\mathbf{Osc}$ (oscillatory, classical
    mechanics language), $\mathbf{Cat}$ (categorical, S-entropy language),
    $\mathbf{Part}$ (partition, quantum mechanics language).
  \item Three explicit conversion functors are constructed:
    $F_{OC} : \mathbf{Osc} \to \mathbf{Cat}$ (frequency divisibility and
    label encoding),
    $F_{CP} : \mathbf{Cat} \to \mathbf{Part}$ (label quotient),
    $F_{PO} : \mathbf{Part} \to \mathbf{Osc}$ (inverse-measure frequency via
    Kac's lemma analogue).
  \item \textbf{Triple Equivalence}: $\mathbf{Osc} \cong \mathbf{Cat} \cong
    \mathbf{Part}$ via the six functors.  The equivalence is exact, not
    approximate.  The symmetry group is $S_3$.
  \item \textbf{Consequence~\ref{cons:classeqquant}}: Classical mechanics
    (trajectories) and quantum mechanics (eigenstates) are equivalent for any
    bounded physical system satisfying Axiom~\ref{ax:bpsl}.
  \item \textbf{Consequence~\ref{cons:gpuphys}}: A GPU fragment shader
    performing partition-depth minimisation on a bounded domain is a physical
    observation apparatus in the formal sense of the calculus.
  \item \textbf{Consequence~\ref{cons:sentunique}}: The three S-entropy
    components $\Sk = \St = \Se$; the representation used for computation is
    a matter of cost, not of physical content.
  \item \textbf{Triple Entropy Equivalence}
    (Consequence~\ref{cons:tripleent}): the oscillatory, categorical, and
    partition entropies are the same Shannon entropy of the partition measure,
    computed in three distinct notations.
  \item Free conversion (Corollary~\ref{cor:free-conv-mono}) allows any
    computation to be relocated to the cheapest representation; apples-and-oranges
    compositions are always admissible in a representation-coherent receiver
    (Corollary~\ref{cor:apples-mono}).
\end{itemize}
\end{tcolorbox}
